\subsection{An explicit output bound for the genus-three construction}
\label{sec:markmanquantitative}

We now carry the numerical calculation through the derived direct image.
The following bounds are new derivations in this run. The geometric inputs
from \citet[Proposition~9.2.1]{Markman25} are the identification of the
output sheaf, its reflexivity and rank, and its codimension-four nonfree
locus. The exact supplied-source locators are
\texttt{assumption-on-C-i-s}, lines~5124--5135;
\texttt{eq-F-1-and-F-2}, lines~5403--5410;
\texttt{prop-local-freeness}, lines~5450--5465; and its proof,
lines~5846--5950. The direct-image bundles and their numerical estimates
below are constructed explicitly.

Let $X=J(C)$ for a smooth nonhyperelliptic genus-three curve. Identify
$\widehat X$ with $X$ by the principal polarization $\Theta$.
Let $Z_1$ be thirteen disjoint translates of the Abel--Jacobi curve, and
let $Z_2$ be thirteen disjoint translates of its reflection. We impose
Markman's general-position hypotheses, including
\begin{equation}\label{eq:M7H0}
 H^0\bigl(X,I_{Z_1}(2\Theta)\otimes P\bigr)=0
 \quad\text{for every }P\in\operatorname{Pic}^0(X).
\end{equation}
The translates may be chosen as the two cyclic orbits required for his
semiregular quotient construction. Put
\[
 A=X_x\times\widehat X_\alpha,\qquad Y=X_y\times A,\qquad
 \pi:Y\longrightarrow A,\qquad a(y,x)=y+x,
\]
and, with the normalized Poincar\'e bundle $\mathcal P$,
\[
 \mathcal L=\cO(\Theta_y)\otimes a^*\cO(\Theta)
                  \otimes\mathcal P_{y,\alpha}^{-1}.
\]
Let $D_1=\{y\in Z_1\}$, $D_2=\{y+x\in Z_2\}$, and set
\[
 \mathcal M=\mathcal L\otimes I_{D_1}I_{D_2},\qquad
 G=R\pi_*\mathcal M,\qquad E=G^\vee[-1].
\]
The sheaf $E$ is reflexive of rank $96$ and is locally free outside a
closed subset of codimension four. Use the very ample class
\[
 H=3(\Theta_x+\Theta_\alpha).
\]
Regularity with respect to $H$ means
$H^i(A,F((m-i)H))=0$ for $i>0$, equivalently regularity of the direct
image under the embedding defined by $H$.

\begin{proposition}[M7: explicit output regularity]\label{M7}
The sheaf $E$ is $112$-regular with respect to $H$. More strongly, if
$m\ge1$ is an integer and $L$ is a line bundle such that
$c_1(L)-112m^2H$ is nef, then
\begin{equation}\label{eq:M7strong}
 H^i(A,[m]^*E\otimes L)=0\qquad(i>0).
\end{equation}
The assertion allows arbitrary numerically trivial factors in $L$.
\end{proposition}

\begin{proof}
\emph{Step 1: an explicit three-term direct image.}
The change of coordinates $(y,x)\mapsto(y,z=y+x)$ identifies the two
curve-union ideals with pullbacks from independent factors. They are Tor
independent, and their product equals their intersection. Consequently
\[
 0\longrightarrow\mathcal M\longrightarrow\mathcal L
 \longrightarrow\mathcal L|_{D_1}\oplus\mathcal L|_{D_2}
 \longrightarrow\mathcal L|_T\longrightarrow0,
 \qquad T=D_1\cap D_2,
\]
is exact on the total space. The scheme $T$ is the disjoint union of
$169$ smooth components $C_i\times\Sigma_j\times\widehat X$ of dimension
five. Its map $f:T\to A$ is finite. Up to translations and a sign, its
first-factor map is
\[
 C\times C\longrightarrow C^{(2)}\simeq\Theta\hookrightarrow X.
\]
Here nonhyperellipticity makes the Abel map $C^{(2)}\to\Theta$ an
isomorphism onto a smooth theta divisor, and the first arrow is finite
flat of degree two.

Define
\[
 W=\pi_*\mathcal L,\qquad
 U=\pi_*(\mathcal L|_{D_1}\oplus\mathcal L|_{D_2}),\qquad
 \mathcal T=f_*(\mathcal L|_T).
\]
On a fibre of $\pi$, $\mathcal L$ has class $2\Theta$, so all its higher
cohomology vanishes and $W$ is locally free of rank eight. Its restriction
to each curve has degree six; since the curve has genus three,
$H^1=0$ and $h^0=4$. Hence $U$ is locally free of rank $26\cdot4=104$.
The finiteness of $f$ removes its higher direct images. Thus $G$ is
represented by
\[
 [\,W\longrightarrow U\longrightarrow\mathcal T\,]
\]
in cohomological degrees $0,1,2$. This uses the exact total-space
sequence; it does not presume left exactness after every fibre
specialization.

By \eqref{eq:M7H0}, $W\to U$ is fibrewise injective: a section in its
kernel vanishes on $Z_1$. Therefore $Q=U/W$ is a vector bundle of rank
$96$, and $G=[Q\to\mathcal T]$ in degrees $1,2$.
Each summand of $\mathcal T$ is a vector bundle on a smooth divisor,
pushed forward to $A$, so it has projective dimension one. Finite
duality gives
\[
 \mathcal T^D:=\mathcal E xt^1_A(\mathcal T,\cO_A)
 =\bigoplus_{i,j}f_{ij,*}
       \bigl((\mathcal L|_{T_{ij}})^{-1}\otimes\omega_{T_{ij}}\bigr),
\]
with $\mathcal H om(\mathcal T,\cO_A)=0$ and higher Ext sheaves zero.
Dualizing $[Q\to\mathcal T]$ and shifting by $-1$ yields
\begin{equation}\label{eq:M7extension}
 0\longrightarrow Q^*\longrightarrow E
          \longrightarrow\mathcal T^D\longrightarrow0.
\end{equation}

\emph{Step 2: the dual bundles.}
For every $q\ge1$,
\begin{equation}\label{eq:M7dualvanishing}
 H^i(A,W^*(qH))=H^i(A,U^*(qH))=0\qquad(i>0).
\end{equation}
For $W$, relative duality gives $W^*=R^3\pi_*\mathcal L^{-1}$.
Leray and Serre duality on the ninefold $Y$ identify the dual of
$H^i(A,W^*(qH))$ with
\[
 H^{6-i}(Y,\mathcal L\otimes\pi^*\cO_A(-qH)).
\]
Project instead onto $X_y$. The line bundle on each $A$-fibre has matrix
$\operatorname{diag}(1-3q,-3q)$ relative to the product principal
polarization. It is antiample, so only the sixth direct image can
be nonzero. Total cohomology below degree six vanishes, as required.

For a curve summand of $U$, relative duality over $A$, followed by Serre
duality on $C\times A$, gives the same group with degree $6-i$ and
line bundle $\mathcal L_C(-qH)$. Projecting onto $C$, the $D_1$ fibre
matrix is again $\operatorname{diag}(1-3q,-3q)$. For $D_2$, use
$z=y+x$ as the fixed curve coordinate. Its matrix is
\[
 \begin{pmatrix}1-3q&1\\1&-3q\end{pmatrix}.
\]
The largest eigenvalue before subtracting $3qI$ is
$(1+\sqrt5)/2<3$, so this is antiample as well. This proves
\eqref{eq:M7dualvanishing}. Translations only change numerically trivial
factors. Since $\dim A=6$, both dual bundles are $7$-regular.

\emph{Step 3: a kernel bound from six Buchsbaum--Rim terms.}
Let $N=h^0(U^*(7H))$. Evaluation and regularity give
\[
 0\longrightarrow K_U\longrightarrow\cO_A(-7H)^N
                  \longrightarrow U^*\longrightarrow0.
\]
The kernel is locally free and $14$-regular. For cohomological degrees
$2\le i\le6$, twist by $(14-i)H$ and use
\eqref{eq:M7dualvanishing} and the positivity of $(7-i)H$.
For $i=1$, the required surjectivity at twist $13H$ follows from
\[
 H^0(U^*(7H))\otimes
       \operatorname{Sym}^6H^0(\mathbf P^M,\cO(1))
       \longrightarrow H^0(U^*(13H)),
\]
which is surjective by $7$-regularity and factors through the evaluation
map with $H^0(A,6H)$.

Compose evaluation with the vector-bundle surjection $U^*\to W^*$,
and call its kernel $J$. There is an exact sequence
\begin{equation}\label{eq:M7kernels}
 0\longrightarrow K_U\longrightarrow J\longrightarrow Q^*
                   \longrightarrow0.
\end{equation}
The Buchsbaum--Rim resolution of $J$ has terms
\begin{equation}\label{eq:M7BR}
 C_k=\bigwedge^{8+k}\bigl(\cO_A(-7H)^N\bigr)
             \otimes\operatorname{Sym}^{k-1}W\otimes\det W,
 \qquad 1\le k\le N-8,
\end{equation}
with $C_1\to J$ surjective. The determinant factor and the indexing are
as in \citet[pp.~41--42]{RaicuSam22}. Exactness here is elementary:
locally split the surjection as $K\oplus V\to V$, $\dim V=8$.
The summand with $a$ factors from $K$ becomes, after
$\bigwedge^{8+k-a}V\otimes\det V^*
\simeq\bigwedge^{a-k}V^*$, the strand
\[
 \bigwedge^{a-k}V^*\otimes\operatorname{Sym}^{k-1}V^*
\]
of the polynomial de Rham complex of total degree $a-1$.
For $a\ge2$ the Euler contracting homotopy proves exactness in
characteristic zero. The $a=1$ strand maps isomorphically onto $K$,
while the remaining $V$ maps isomorphically onto the target.
Thus the augmented complex is locally split exact.

The theorem of the square, with $b(x,\alpha)=x-\alpha$, gives
$W=\cO_A(\Theta_x)\otimes b^*\mathrm{FM}(\cO_X(2\Theta))$.
The pullback formula for this Fourier--Mukai bundle gives
\begin{equation}\label{eq:M7splitW}
 [2]^*W\simeq Q_0^{\oplus8},\qquad
 c_1(Q_0)=4\Theta_x-2\Theta_{x-\alpha}.
\end{equation}
To check the formula directly, for $L=\cO_X(2\Theta)$ pull the
Poincar\'e family back by its polarization isogeny $\phi_L=[2]$.
The theorem of the square identifies it with
$a^*L\otimes p_2^*L^{-1}$, whose direct image is
$H^0(L)\otimes L^{-1}$. Changing the Poincar\'e sign changes $[2]$
to $[-2]$ and leaves the numerical class unchanged.
Different symmetric normalizations introduce only a numerically trivial
factor, harmless for every ampleness test below.

The matrix of $c_1(Q_0)$ is
$\left(\begin{smallmatrix}2&2\\2&-2\end{smallmatrix}\right)$,
with eigenvalues $\pm2\sqrt2$, whereas $[2]^*H$ has matrix $12I_2$.
For $q\ge k+7$, the line bundle with class
$q[2]^*H+(k+7)c_1(Q_0)$ is ample.
By \eqref{eq:M7splitW}, Kodaira vanishing and the trace splitting
for $[2]$,
\[
 \operatorname{reg}_H
   (\operatorname{Sym}^{k-1}W\otimes\det W)\le k+13.
\]
It follows from \eqref{eq:M7BR} that
\[
 \operatorname{reg}_H C_k\le7(8+k)+k+13=69+8k.
\]
Only six terms are needed: successive cohomology sequences show that
$H^i(J((R-i)H))=0$ if
$H^{i+k-1}(C_k((R-i)H))=0$ for $1\le k\le7-i$.
The final kernel then has cohomology in degree seven, which vanishes
by dimension. Therefore
\[
 \operatorname{reg}_H J
 \le\max_{1\le k\le6}\{\operatorname{reg}_H C_k-k+1\}
 \le\max_{1\le k\le6}(70+7k)=112.
\]
Sequence \eqref{eq:M7kernels} and $14$-regularity of $K_U$ give
$112$-regularity of $Q^*$. No numerical bound for $N$ was needed.

\emph{Step 4: the torsion dual.}
On a component $T=C\times C\times\widehat X$, let
$F_1=\{p\}\times C$, $F_2=C\times\{p\}$,
$S=F_1+F_2$, and let $\Delta_C$ denote the diagonal.
The addition map to $X$ pulls $\Theta$ back to
\[
 A_C=4S-\Delta_C.
\]
Indeed $\Theta|_C$ has degree three and the mixed Poincar\'e class
is $S-\Delta_C$; the formula also gives degree twelve on the diagonal,
as required by $[2]^*\Theta|_C$.
The class $R_C=3S-\Delta_C$ is nef. Choose a degree-three morphism
$C\to\mathbf P^1$ by projection of the plane quartic from one of its
points. The residual divisor to $\Delta_C$ in its fibre correspondence
has class $R_C$. Its transitive monodromy is $S_3$ or $A_3$.
In the $S_3$ case the residual divisor is irreducible, since ordered
unequal pairs form a transitive orbit. Its self-intersection is
\[
 R_C^2=9S^2-6S\Delta_C+\Delta_C^2=18-12-4=2.
\]
An effective irreducible divisor of nonnegative self-intersection on a
smooth surface is nef. In the $A_3$ case the cover is cyclic, and the
residual divisor is $\Gamma_\sigma+\Gamma_{\sigma^2}$.
Each graph has self-intersection $-4$; Riemann--Hurwitz gives five
fixed points, at which the graphs meet transversely.
Thus its intersection with either graph is $-4+5=1$, and its
intersection with every other irreducible curve is nonnegative.
This proves nefness in both cases.

Numerically,
\[
 c_1(\mathcal L_T)=3F_1+3F_2-\mathcal P_{c_1,\alpha},
 \qquad f^*H=3A_C+3\Theta_\alpha.
\]
For $q\ge3$ write
\begin{align*}
 qf^*H-c_1(\mathcal L_T)
  ={}&[\,3qA_C-6F_1-3F_2\,]+(3q-1)\Theta_\alpha\\
    &+[\,3F_1+\Theta_\alpha+\mathcal P_{c_1,\alpha}\,].
\end{align*}
The last bracket is a pullback of $\Theta$ under
$(c_1,\alpha)\mapsto AJ(c_1)+\alpha$, hence nef.
The first is
$3qR_C+(3q-6)F_1+(3q-3)F_2$, which is ample on $C\times C$.
Together with the positive theta class on $\widehat X$ it is ample
on the product, and adding the last nef class preserves ampleness.
Kodaira vanishing on the smooth fivefold $T$ yields
\[
 H^i(T,\omega_T\otimes\mathcal L_T^{-1}\otimes f^*\cO_A(qH))=0
 \quad(i>0,\ q\ge3).
\]
Thus $\mathcal T^D$ is $8$-regular: its support has dimension five,
and $8-i\ge3$ for $1\le i\le5$.
The extension \eqref{eq:M7extension} now proves
$112$-regularity of $E$.

\emph{Step 5: the stronger estimate.}
Pull all the displayed exact complexes back by the finite \'etale
map $[m]$. The duality argument of Step~2 applies after this base
change: the two fibre matrices are multiplied by $m^2$.
Subtracting $c_1(L)$ makes them negative definite if
$c_1(L)-112m^2H$ is nef. Hence $[m]^*U^*\otimes L$ and
$[m]^*W^*\otimes L$ have no higher cohomology.
The pulled-back evaluation sequence has middle term
$L(-7m^2H)^N$, an ample line bundle sum. It follows that
\[
 H^j(A,[m]^*K_U\otimes L)=0\qquad(j\ge2).
\]
No assertion about its $H^1$, and no multiplication map with a
varying flat factor, is needed.

For $k\le6$, pull $[m]^*C_k\otimes L$ back once more by $[2]$.
Its summands have class
\[
 [2]^*c_1(L)-7m^2(8+k)[2]^*H+m^2(k+7)c_1(Q_0).
\]
After removing a nef class, the matrix is $m^2$ times
\[
 12[112-7(8+k)]I_2+
       (k+7)\begin{pmatrix}2&2\\2&-2\end{pmatrix}.
\]
For $1\le k\le6$ its least eigenvalue is positive; the smallest
lower bound used here is $168-26\sqrt2>0$.
Kodaira vanishing and trace therefore give higher-cohomology
vanishing for the first six terms. The same truncation argument,
then \eqref{eq:M7kernels}, proves it for $[m]^*Q^*\otimes L$.

For the final term let $T'=T\times_{A,[m]}A$, with
$g:T'\to T$ finite \'etale and $f':T'\to A$ finite.
Finite flat base change gives
$[m]^*\mathcal T^D=f'_*(\omega_{T'}\otimes g^*\mathcal L_T^{-1})$.
The class
\[
 f'^*c_1(L)-g^*c_1(\mathcal L_T)
 =g^*(3f^*H-c_1(\mathcal L_T))
       +f'^*(c_1(L)-3m^2H)
\]
is ample, by Step~4 and the nef hypothesis.
Kodaira vanishing on $T'$ and
\eqref{eq:M7extension} finish the proof of \eqref{eq:M7strong}.
\end{proof}

\subsection{Effective Chern cycles and the normalized family}

\begin{proposition}[M8: effective third-Chern representatives]\label{M8}
Let $B$ be a rank-$r$ coherent sheaf, $r\ge3$, on a smooth projective sixfold $A$,
locally free away from a closed subset of codimension at least four.
Let $h$ be very ample, let $R\ge0$, and suppose $F=B(Rh)$ is globally
generated. Assume $c_1(B)=0$ in rational cohomology. Then $c_3(B)$ has
an integral signed effective representative $Z^+-Z^-$ with
\begin{align}
 D^+&:=\deg_h Z^+
   =\int h^3c_3(B)+(r-2)R\int h^4c_2(B)
       +\frac{r(r-1)(r-2)}2R^3\int h^6,\label{eq:M8degreesplus}\\
 D^-&:=\deg_h Z^-
   =(r-2)R\int h^4c_2(B)
       +\frac{r(r-1)(r-2)}2R^3\int h^6.\label{eq:M8degreesminus}
\end{align}
The cycle identities hold in Chow groups if $c_1(B)=0$ there; the
stated cohomological assumption suffices for the cycle-class and
degree conclusions.
\end{proposition}
\begin{proof}
On the locally free open set choose $r-p+1$ general sections of $F$,
for $p=2,3$. Global generation makes their universal evaluation
transverse to the matrix-rank strata. Generic smoothness gives an
expected codimension-$p$ dependence locus, or the empty locus.
Its class is $c_p(F)$, by the elementary maximal-rank-drop case of
the Thom--Porteous formula. Let $C_p$ be its effective closure.

If $S$ is the nonfree locus, then $\dim S\le2$.
The localization sequence gives
$\CH^p(A)\simeq\CH^p(A\setminus S)$ for $p\le3$, because
$S$ has no cycle of dimension $6-p$. Consequently $C_p$ represents
$c_p(F)$ globally. This also explains why possible excess loci
inside $S$ contribute no unwanted codimension-two or codimension-three
terms.

The twist identities, after passing to the stated cohomology if
necessary, are
\[
 c_2(F)=c_2(B)+\binom r2R^2h^2,\qquad
 c_3(F)=c_3(B)+(r-2)Rh\,c_2(B)+\binom r3R^3h^3.
\]
Eliminating $c_2(B)$ gives
\begin{equation}\label{eq:M8cycleidentity}
 c_3(B)=C_3-(r-2)Rh\,C_2+
                \frac{r(r-1)(r-2)}3R^3h^3.
\end{equation}
The last coefficient is an integer. General divisors in $|h|$ meet
$C_2$ properly, and three such divisors represent $h^3$ effectively.
Thus take
\[
 Z^+=C_3+\frac{r(r-1)(r-2)}3R^3h^3,\qquad
 Z^-=(r-2)Rh\,C_2.
\]
Intersecting with $h^3$ and expanding $c_2(F),c_3(F)$ gives
\eqref{eq:M8degreesplus}--\eqref{eq:M8degreesminus}.
Both displayed integers are nonnegative because the constructed
cycles are effective.
\end{proof}

Apply this to the output $E$ of M7. Let $r=96$, $D=\det E$, and let
$h$ be a symmetric very ample line bundle extending as the fixed
polarization in the Weil deformation family. Suppose integers $c,C>0$
satisfy
\begin{equation}\label{eq:M8comparisons}
 ch-H\ \text{ample},\qquad Ch-c_1(D)\ \text{ample}.
\end{equation}
Here $ch$ means the scalar multiple $c\,h$, not the Chern character.
Define the ordinary normalized seed
\[
 B_0=[r]^*E\otimes D^{-r}.
\]
Its first Chern class vanishes in cohomology and
\begin{equation}\label{eq:M8normalizedcharacter}
 \ch_p(B_0)=r^{2p}\kappa_p(E),\qquad
 c_3(B_0)=2r^6\kappa_3(E).
\end{equation}
Indeed $[r]^*c_1(E)=r^2c_1(E)$, while the rank-$r$ tensor correction
subtracts $r^2c_1(E)$, and the same computation in the exponential
Chern character proves the first identity.

For seed global generation one can already use
$R_0=112r^2c+rC$: $E(112H)$ is globally generated, so tensoring by
the ample power $112(ch-H)$ gives $E(112ch)$ globally generated.
Pull back by $[r]$, then tensor by the globally generated line bundle
$(h^C\otimes D^{-1})^r$.
Squares of ample line bundles on abelian varieties are globally
generated, including arbitrary flat twists
\citep[Theorem~1.1(a)]{BauerSzemberg96}.
The determinant of $B_0$ can have a numerically trivial factor.
Any compatible determinant normalization changes $B_0$ only by a
flat line bundle and is covered by the same argument.

For persistence in a family, global generation at the seed alone
would not suffice. M7 provides the stronger explicit regularity bound
\begin{equation}\label{eq:M8R}
 R_{\mathrm{safe}}=112r^2c+rC+6.
\end{equation}
For $1\le i\le6$ take
$L_i=D^{-r}\otimes h^{R_{\mathrm{safe}}-i}$, including any required
flat normalization factor. Then
\[
 c_1(L_i)-112r^2H
 =112r^2(ch-H)+r(Ch-c_1(D))+(6-i)h
\]
is ample. Equation \eqref{eq:M7strong}, with $m=r$, proves all six
vanishings
\[
 H^i(A,B_0((R_{\mathrm{safe}}-i)h))=0.
\]
Thus the normalized seed is $R_{\mathrm{safe}}$-regular.

For precision about the family step, use the flat determinant-normalized
$\mu_r$-twisted deformation supplied by Markman's semiregularity
construction, pulled back from his quotient to the original torus
family. On a sufficiently small contractible analytic base,
$[r]^*$ kills its twisting class: multiplication acts as $r^2$ on
the local system $H^2(-,\mu_r)$. A $\mu_r$-\v Cech trivialization
over this smaller family produces an ordinary flat coherent family.
Its seed is the displayed $B_0$ up to a flat line bundle.
Analytic upper semicontinuity
\citep[\S7, Satz~3, pp.~62--63]{Grauert60}
preserves the six proved vanishings on a neighbourhood. Its
hypotheses hold because the torus family is proper over a reduced
analytic base and the untwisted coherent sheaf is flat over it.
The locus where the flat family is locally free is
open in the total space; properness gives upper semicontinuity of
the dimension of its closed nonfree fibres. Its codimension is
therefore at least four on a smaller neighbourhood.
M8 applies fibrewise with the same $R_{\mathrm{safe}}$ and the same
intersection numbers. GAGA makes these analytic-fibre sheaves and
their cycles algebraic. The bounded-Chow argument R21 will be used
after the fixed lattice and projector conversions below.

If the Weil projection of $c_3(B_0)$ is nonzero, at least one of
$Z^+,Z^-$ has nonzero Weil projection. This supplies exactly the
degree-only input of P3. No estimate for its direction relative to
a prescribed Weil vector is needed at this stage.

\subsection{Exact GRR and a numerical source degree}

\begin{proposition}[M9: a concrete effective sixfold seed]\label{M9}
For the genus-three $d=12$ construction, take
$h=3(12\Theta_x+\Theta_\alpha)$. On a nonempty analytic open subset
of its nine-dimensional Weil period family there is an integral
effective three-cycle with nonzero Weil projection and degree at most
\begin{equation}\label{eq:M9Dsource}
 D_{\mathrm{src}}=
 3\,426\,088\,289\,770\,714\,707\,027\,306\,735\,206\,400.
\end{equation}
The normalized ordinary sheaf used to construct it is
$2\,065\,542$-regular with respect to $h$.
\end{proposition}
\begin{proof}
Put $x=\Theta_x$, $y=\Theta_\alpha$, and
$p=c_1(\mathcal P_{x,\alpha})$. Our matrix convention is
\[
 bx+cy+ep\ \longleftrightarrow\
 \begin{pmatrix}b&e\\e&c\end{pmatrix},\qquad
 \int_{X\times\widehat X}\exp(bx+cy+ep)=(bc-e^2)^3.
\]
The input character is
$\ch(F)=1+\theta-6\theta^2-2\theta^3$, from the displayed secant
character in Markman's supplied source, lines~3743--3776.
Set $(a_0,a_1,a_2,a_3)=(1,1,-6,-2)$ and let
\[
 \mathcal D=\left.
       \sum_{i,j=0}^3 a_i a_j\partial_u^i\partial_v^j
       \right|_{u=v=0}.
\]
Factoring exterior integration into the three symplectic coordinate
pairs gives
\[
 \pi_*\exp(u\theta_y+v\theta_{y+x}-\mathcal P_{y,\alpha})
 =(u+v)^3
       \exp\left(\frac{uvx+vp-y}{u+v}\right).
\]
For $k\le3$ the degree-$2k$ part is the polynomial
$(u+v)^{3-k}(uvx+vp-y)^k/k!$. Apply $\mathcal D$ and multiply
by $(-1)^{k+1}$, since $E=G^\vee[-1]$. One obtains
\begin{align*}
 \ch_0(E)&=96,\\
 \ch_1(E)&=264x+22y-48p,\\
 \ch_2(E)&=12p^2-132px-11py-144x^2-24xy-y^2,\\
 \ch_3(E)&=-2p^3-6p^2x+6p^2y+72px^2+12pxy+\tfrac12py^2\\
          &\hspace{9mm}+24x^3-72x^2y+\tfrac12xy^2-\tfrac16y^3.
\end{align*}
Consequently
\begin{align}
 \kappa_2(E)&=-\frac{169}{48}(12x+y)^2,\label{eq:M9kappa2}\\
 \kappa_3(E)&=\frac{13}{3456}
 \bigl(-10368p^2x+864p^2y+288576x^3\nonumber\\
 &\hspace{27mm}+51408x^2y+6012xy^2+119y^3\bigr).
 \label{eq:M9kappa3}
\end{align}

These formulas also identify the invariant polarization without
guessing a dual-torus sign convention. Markman's normalized character
is $\operatorname{Spin}(V)_P$-invariant, as in his main theorem,
item~3, and supplied-source lines~2406--2468.
Let $h_0=12x+y$. Differentiating the invariance of the nonzero
multiple $h_0^2$ in \eqref{eq:M9kappa2} gives
$h_0\wedge(a\cdot h_0)=0$ for every Lie-algebra element $a$.
Since $h_0$ is ample, Lefschetz injectivity $H^2\to H^4$ gives
$a\cdot h_0=0$. The connected group fixes $h_0$.
Its invariant line in $H^2$ is one-dimensional
(source lines~3966--3968), so $h_0$ is the required integral
polarization. Take $h=3h_0$. The third power of every ample line bundle
on an abelian variety is very ample
\citep[Theorem~1.1(b)]{BauerSzemberg96}; a local relative symmetric
representative therefore gives the required relative embedding
throughout the polarized deformation.

At the seed the polarization type of $h$ is
$(3,3,3,36,36,36)$, and hence
\[
 \Delta_h:=\frac{h^6}{6!}=108^3=1\,259\,712,\qquad
 \int h^6=906\,992\,640.
\]
The determinant class of $E$ has matrix
$\left(\begin{smallmatrix}264&-48\\-48&22\end{smallmatrix}\right)$.
Choose $c=2$, $C=12$. The two comparison matrices in
\eqref{eq:M8comparisons} are
\[
 2h-H=\begin{pmatrix}69&0\\0&3\end{pmatrix},\qquad
 12h-c_1(\det E)=\begin{pmatrix}168&48\\48&14\end{pmatrix}.
\]
Their leading principal minors are $(69,207)$ and $(168,48)$,
so both classes are ample. Equation \eqref{eq:M8R} gives
\[
 R=112\cdot96^2\cdot2+96\cdot12+6=2\,065\,542.
\]

For $B_0=[96]^*E\otimes(\det E)^{-96}$, use
\eqref{eq:M8normalizedcharacter} and the intersection polynomial
above to obtain
\begin{align*}
 \int h^4c_2(B_0)&=30\,136\,419\,515\,105\,280,\\
 \int h^3c_3(B_0)&=28\,288\,052\,451\,512\,156\,160.
\end{align*}
For example $c_2(B_0)=33\,226\,752h^2$ in rational cohomology.
The coefficient of its $h^3$ component in $c_3(B_0)$ is
$31\,188\,844\,544$; the remaining Weil component is nonzero by
Markman's independence-from-$h^3$ assertion.
Substitution into M8, with $r=96$, gives
\begin{align*}
 D^-&=94R\int h^4c_2(B_0)+428640R^3\int h^6\\
    &=3\,426\,088\,289\,770\,686\,418\,974\,855\,223\,050\,240,\\
 D^+&=D^-+\int h^3c_3(B_0)\\
    &=3\,426\,088\,289\,770\,714\,707\,027\,306\,735\,206\,400.
\end{align*}
At least one sign has nonzero Weil projection.
The regularity and nonfree-locus arguments preceding M9 preserve
these bounds on an analytic neighbourhood, proving the proposition.

The exact calculations were executed in
\path{work/markman_output_numbers.py}.
As a distinct integration-order check, the same operator
$\mathcal D$ was applied to the cube of the determinant of
\[
 \begin{pmatrix}
 u+v&v&-1\\
 v&v+b&e\\
 -1&e&c
 \end{pmatrix}.
\]
This is absolute GRR on $X^3$. Its top four homogeneous components
agree with the four characters just displayed paired against
$\exp(bx+cy+ep)$. This verifies the arithmetic signs and factorials;
the geometric proof is the argument above.
\end{proof}

\subsection{Maximal-order transfer and an explicit fourfold cap}

We first transfer the analytic-open seed to the principal split
sixfold family. This keeps the subsequent integral lattice calculation
in a fixed, completely specified lattice.

\begin{lemma}[M10a: principal-sixfold transfer]\label{M10a}
Set
\[
 M=4\cdot16^6\cdot108^3,\qquad
 D_6=16^3M^6D_{\mathrm{src}}.
\]
On a nonempty analytic open subset of the principal split sixfold
family over $\Q(\omega)$ there is an integral effective three-cycle
of principal-polarization degree at most $D_6$ with nonzero Weil
projection.
\end{lemma}
\begin{proof}
At the pure-spinor seed, $f=\sqrt{-12}$ is integral on the homology
lattice $\Lambda$. In fact the source eigenspaces
\texttt{eq-W-1-and-2}, lines~1353--1364, give its cohomological
block matrix
\[
 \begin{pmatrix}0&12\theta\\-\theta^{-1}&0\end{pmatrix},
\]
where $\theta$ is the unimodular principal-polarization contraction.
The dual matrix is integral on homology. Since
$\omega=(f-2)/4$, the lattice
\[
 \Lambda_{\cO}=\Lambda+\omega\Lambda
 \subset\tfrac14\Lambda
\]
is $\cO_K$-stable, using $\omega^2+\omega+1=0$.
Its index $I=[\Lambda_{\cO}:\Lambda]$ divides $4^{12}$.
The identity on the universal cover defines a quotient isogeny
$q_{\cO}:A\to A_{\cO}$ of degree $I$.

The Riemann form $16h$ is integral on $\Lambda_{\cO}$: each argument
has denominator dividing four. It therefore defines an ample
polarization $h_{\cO}$ with
$q_{\cO}^*h_{\cO}=16h$. This equality can be arranged for line
bundles as well, since isogeny pullback is surjective on
$\operatorname{Pic}^0$. Its polarization product is
\[
 \Delta_{\cO}=\frac{16^6\Delta_h}{I}\in\Z,
 \qquad \Delta_h=108^3.
\]
Thus $\Delta_{\cO}$ is an integer divisor of $16^6\Delta_h$.
All these lattice constructions are constant on a marked analytic
neighbourhood.

Here is the required quantitative Hermitian comparison.
For an $\cO_K$-stable lattice with compatible alternating form
$\mathcal E$, use
\[
 \mathfrak h(s,t)=\mathcal E(s,\omega t)
                         -\bar\omega\,\mathcal E(s,t).
\]
It is integral Hermitian, conjugate-linear in the first variable;
the coefficient of $\omega$ recovers $\mathcal E$.
Its absolute determinant equals the product $\Delta$ of the
polarization elementary divisors: on the underlying integral lattice
the alternating determinant is $\Delta^2$.

In a split Hermitian space of $K$-dimension six, intersect a rational
maximal isotropic three-plane with the lattice. The intersection is
a primitive $\cO_K$-submodule. Since $\Z[\omega]$ is a PID, its basis
$v$ extends to a lattice basis $(v,w)$, with Gram matrix
\[
 \begin{pmatrix}0&B\\B^*&D\end{pmatrix},
 \qquad |\det\mathfrak h|=\Nm(\det B)=\Delta.
\]
In this odd half-rank case the determinant itself is negative; its
absolute value is used throughout. Put
\[
 X=(B^*)^{-1}D/2,\qquad w'=w-vX,\qquad u=w'B^{-1}.
\]
Then $(v,u)$ has hyperbolic Gram matrix. If $b=\det B$, the entries
of $X$, $B^{-1}$ and $XB^{-1}$ have denominators dividing
$2\bar b$, $b$, and $2b\bar b=2\Delta$, respectively.
The isometry to a hyperbolic basis and its inverse are explicitly
\[
 \mu(v)=e,\quad \mu(w)=eX+fB,\qquad
 \mu^{-1}(e)=v,\quad \mu^{-1}(f)=wB^{-1}-vXB^{-1}.
\]
Thus both directions have denominator dividing $2\Delta$.
Apply this to $A_{\cO}$ and to the fixed principal split target
lattice, whose product is one. Their comparison has denominator
dividing $4\Delta_{\cO}$, hence dividing $M$.

Transport the rational comparison across the marked period
neighbourhood. Multiplication by $M$ gives a $K$-linear isogeny
$\phi:A_{\cO}\to A_{\mathrm{pr}}$ satisfying
$\phi^*\theta_{\mathrm{pr}}=M^2h_{\cO}$.
The rational isometry identifies the period domains, so the image
of the source analytic open set is again analytic open.
For an effective three-cycle $Z$, the projection formula gives
\[
 \deg_{\theta_{\mathrm{pr}}}
       ((\phi\circ q_{\cO})_*Z)=16^3M^6\deg_h Z.
\]
Isogeny pushforward acts invertibly on rational Weil spaces, so it
preserves a nonzero Weil projection. M9 proves the claim.
Only ampleness of $h_{\cO}$ is needed. If a projective embedding of
the principal target is required for R21, use $3\theta_{\mathrm{pr}}$
and the fixed degree factor $3^3$; the intrinsic bound remains
measured by $\theta_{\mathrm{pr}}$.
\end{proof}

For the next statement choose the saturated integral Weil basis
$u_3,v_3$ on the principal sixfold component using the balanced
$E^6$ marking, and pass to the finite determinant cover so that
these classes are flat. Its positive Gram matrix is
\[
 G_3=9\begin{pmatrix}2&1\\1&2\end{pmatrix}.
\]
For completeness, write $e_j,f_j$ for the integral basis on each
elliptic factor with $\int e_jf_j=1$, and put
\[
 \Omega_3=\bigwedge_{i=1}^3(e_{x_i}+\omega f_{x_i})
          \wedge\bigwedge_{i=1}^3(e_{y_i}+\omega^2f_{y_i})
          =u_3+\omega v_3.
\]
The coefficients of the all-$e$ monomial and of the monomial
obtained by replacing only $e_{x_1}$ by $f_{x_1}$ are $1$ and
$\omega$. They show that $u_3,v_3$ span the saturated integral
Weil lattice, since $\Omega_3,\bar\Omega_3$ span its two complex
eigenlines. Furthermore $\Omega_3^2=\bar\Omega_3^2=0$ and
$\int\Omega_3\bar\Omega_3=-27$. Expanding gives
$\int u_3^2=\int v_3^2=-18$ and $\int u_3v_3=-9$;
the positive Hodge--Riemann form is minus this intersection form.
These equalities also fix the orientation used in the
compensation calculation below.
For the projector P1, use $a=3+\omega$, and write
$Q_3=\delta_3P_3^{\mathrm{proj}}$ to distinguish this polynomial
from the seed criterion P3. The explicit values are
\begin{align*}
 \delta_3&=1\,701\,760\,099\,946\,011\,389\,900,\\
 B_3&=6\,493\,952\,361\,560\,098\,735\,300,\\
 a^6&=37+360\omega.
\end{align*}
For reproducibility, the coefficients of $Q_3(T)$ in increasing
degree are
\[
\begin{split}
(&-266808926179074658393,\ 2660867472861128286,\\
 &-13456515557549106,\ 43596368737922,\ -93608421858,\\
 &121120686,\ -62245).
\end{split}
\]
They give $Q_3(a^{6-j}\bar a^j)=\delta_3$ for $j=0,6$ and zero
for $1\le j\le5$; these seven values were independently checked
by exact arithmetic in $\Z[\omega]$. The seven eigenvalues are
distinct, and
$B_3=\sum_j |[T^j]Q_3|\,7^{3j}$.
Define the completely explicit integer
\begin{equation}\label{eq:M10H3}
 H_3=\max\left\{1,
   \left\lfloor\frac{\delta_3^2D_6^2}{648}\right\rfloor\right\},
 \qquad L_3=\operatorname{lcm}(1,2,\ldots,H_3).
\end{equation}

\begin{proposition}[M10b: sixfold multiplier and degrees]\label{M10b}
Every fibre of the principal split sixfold component has integral
effective cycles $Z_s^\pm$ satisfying
\[
 N_su_3=[Z_s^+]-[Z_s^-],\qquad
 1\le N_s\le360H_3,\qquad
 \deg Z_s^\pm\le3N_sB_3D_6.
\]
A single explicit multiplier valid on every fibre is
$N_3=360L_3$, with each sign of degree at most $3N_3B_3D_6$.
\end{proposition}
\begin{proof}
On the analytic open subset of M10a, choose an effective cycle
$Z$ with degree at most $D_6$ and nonzero Weil projection.
Put
$z=Q_3(a^*)[Z]=xu_3+yv_3$.
This is a nonzero integral Weil class, and the polynomial gives
a signed representative with each sign of degree at most $B_3D_6$.
The orthogonality and comass argument of P3, with principal volume
one and $\operatorname{vol}(Z)=\deg Z/3!$, gives
\[
 \|z\|\le\frac{\delta_3D_6}{6},\qquad
 \|z\|^2=18(x^2+xy+y^2).
\]
Thus the positive integer $h_z=x^2+xy+y^2$ is at most $H_3$.

Pullback by $a$ acts on this basis by
\[
 T=\begin{pmatrix}37&360\\-360&-323\end{pmatrix},\qquad
 \det(z,Tz)=-360h_z.
\]
Cramer's rule therefore represents
$N_zu_3$, $N_z=360h_z$, by the two integral classes $z,Tz$.
The absolute values of its coefficients are
$|360x+323y|$ and $|y|$.
The cycle-degree multiplier of $a^*$ is $7^3=343$, so its cost
is bounded by
\[
 \bigl(|360x+323y|+343|y|\bigr)B_3D_6.
\]
For integral $x,y$ one has $|x|,|y|\le x^2+xy+y^2=h_z$.
For example the quadratic form is at least $3x^2/4$; for
$|x|\ge2$ this is at least $|x|$, while the remaining cases are
immediate. Hence the cost is at most
$1026h_zB_3D_6<3N_zB_3D_6$.

For each of the finitely many possible $h_z\in[1,H_3]$, use the
closed locus of R21 for multiplier $360h_z$ and degree bound
$3(360h_z)B_3D_6$. Their union contains a nonempty analytic open
subset, so irreducibility makes it the entire component.
This preserves the stated relation between each multiplier and
its degree bound. Alternatively each $N_z$ divides $N_3=360L_3$.
Scale the open-set representatives first; their degrees are then
at most $3N_3B_3D_6$. R21 applied to this fixed multiplier proves
the single-multiplier assertion.
\end{proof}

\begin{theorem}[M10/T5: an explicit fourfold template]\label{M10}
Let $A_4$ vary in the principal split fourfold component for
$K=\Q(\omega)$, and let $u_2,v_2$ be its saturated integral flat
Weil basis in the balanced marking. For
\[
 t=\frac{a u_2+b v_2}{q},\qquad a,b\in\Z,\quad q\ge1,
\]
there are integral effective codimension-two cycles satisfying
\[
 Nt=[Z^+]-[Z^-],\qquad
 N\le1080qH_3,\qquad
 \deg Z^\pm\le3240H_3B_3D_6(|a|+|b|).
\]
A single explicit multiplier and degree bound for every fibre are
\[
 N=1080qL_3,\qquad
 \deg Z^\pm\le3240L_3B_3D_6(|a|+|b|).
\]
Here $D_6$, $H_3$, and $L_3$ are the finite formulas in M10a and
\eqref{eq:M10H3}, using the integer \eqref{eq:M9Dsource}.
\end{theorem}
\begin{proof}
Fix $B=E_x\times E_y$ with balanced $K$-action and product principal
polarization $\theta_B$. The compensated variety
$A_6=A_4\times B$ belongs to the principal split sixfold family,
with $\theta_6=\theta_4+\theta_B$.
On $B$ the integral Weil divisors from P2b are
\[
 U=\Gamma_\omega+\Gamma_1-2F_x-2F_y,\qquad
 V=\Gamma_1-F_x-F_y,
\]
and
$\int U^2=\int V^2=-2$, $\int UV=-1$.
With the marked exterior forms,
$\Omega_3=\Omega_2\Omega_1$ and therefore
$u_3=u_2U-v_2V$. Direct pairing yields
\[
 \operatorname{pr}_{A_4,*}\bigl(u_3(-2U+V)\bigr)=3u_2,
 \qquad
 \operatorname{pr}_{A_4,*}\bigl(u_3(-U+2V)\bigr)=3v_2.
\]

The signed divisors in these contractions are concretely
\[
 -2U+V=3F_x+3F_y-2\Gamma_\omega-\Gamma_1,\qquad
 -U+2V=\Gamma_1-\Gamma_\omega.
\]
Each graph divisor is nef and bounded above by $2\theta_B$ in the
nef cone. Thus the first divisor's positive and negative parts
are bounded above by $3\theta_B$ and $6\theta_B$, respectively;
the second's two parts are each bounded by $2\theta_B$.
Translate these finitely many divisors generically on $B$ before
intersecting with any fixed cycle. This makes all intersections
proper, including for cycles with zero-dimensional image on $B$,
and preserves their cohomology classes and numerical bounds.

If $Z$ is an effective three-cycle on $A_4\times B$, nef
intersection positivity gives
\[
 3\int_Z\theta_4^2\theta_B
 \le\int_Z(\theta_4+\theta_B)^3.
\]
Consequently a signed three-cycle whose two signs each have degree
at most $D'$ gives, after either signed divisor contraction, two
effective signs of degree at most $3D'$.
For the first divisor the two separate contributions are bounded
by $(3+6)D'/3$; for the second they are bounded by $4D'/3$.
Proper pushforward preserves effectivity and these degree
identities; components of smaller image dimension contribute zero.

Use the cycle for $N_su_3$ from M10b. The two contractions give
representatives of both $3N_su_2$ and $3N_sv_2$, with the same
multiplier and each sign of degree at most $9N_sB_3D_6$.
Taking their integral combination with coefficients $a,b$ gives
the class $3qN_st$, with each degree at most
$9N_sB_3D_6(|a|+|b|)$. Since $N_s\le360H_3$, this is the first
assertion. Using the single $N_3=360L_3$ gives the second.
\end{proof}

The stated constants concern the principal split datum. A different
fixed integral polarized $K$-datum can be compared using the
Hermitian-lattice calculation above, with its polarization product
and the corresponding fixed rational target-coordinate factors
retained explicitly. Neither these factors nor any raw Hecke
denominator are silently absorbed into M10.

\paragraph{Verification scope.}
M7.1--M7.5 and M8 received a complete independent second proof check
by the reductions agent, recorded as R28.1--R28.6 in
\path{work/markman_regularity_independent_check.md}.
The complete M9--M10 proofs and the final text received a further
independent check, recorded as R30.1--R32.3 in
\path{work/final_integration_review.md} and
\path{work/reductions_check_events.tsv}.
The arithmetic script
\path{work/markman_output_numbers.py} records the two GRR
integration orders, both positivity matrices, the Chern intersection
numbers, the source degree, the maximal-order transfer constants,
and an independent $\Z[\omega]$ Horner evaluation of all seven
projector eigenvalues. It outputs
\path{work/markman_output_numbers.json}; it does not expand the
huge least common multiple. A separate implementation using exact
rational multinomial sums,
\path{work/check_markman_numbers_independent.py}, agrees with
every M9 character coefficient, intersection number, and numerical
M10 cap. These are exact ordinary computations
and targeted mathematical proof checks, not formal
proof-assistant certificates.
